\documentclass[a4paper,12pt]{article}
\usepackage{amsmath,amssymb,amsthm}
\usepackage[margin=1in]{geometry}

\newtheorem{theorem}{Theorem}
\newtheorem{lemma}[theorem]{Lemma}

\title{Problem 342: The Totient of a Square Is a Cube}
\author{Project Euler}
\date{}

\begin{document}
\maketitle

\section{Problem Statement}
Find the sum of all $n$ with $1 < n < 10^{10}$ such that $\varphi(n^2)$ is a perfect cube.

\section{Mathematical Foundation}

\begin{theorem}[Totient of a Square]
For $n = \prod_{i=1}^{k} p_i^{a_i}$,
\[
\varphi(n^2) = \prod_{i=1}^{k} p_i^{2a_i - 1}(p_i - 1).
\]
\end{theorem}

\begin{proof}
By multiplicativity, $\varphi(n^2) = \prod \varphi(p_i^{2a_i}) = \prod p_i^{2a_i - 1}(p_i - 1)$.
\end{proof}

\begin{theorem}[Cube Condition]
$\varphi(n^2)$ is a perfect cube if and only if for every prime $q$,
\[
v_q\!\left(\varphi(n^2)\right) = \sum_{\substack{i:\, p_i = q}} (2a_i - 1) + \sum_{i=1}^{k} v_q(p_i - 1) \equiv 0 \pmod{3},
\]
where $v_q(m)$ denotes the $q$-adic valuation of $m$.
\end{theorem}

\begin{proof}
A positive integer is a perfect cube iff every prime appears in its factorization with exponent divisible by~3. The exponent of $q$ in $\varphi(n^2) = \prod p_i^{2a_i-1}(p_i - 1)$ is exactly the expression above.
\end{proof}

\begin{lemma}[Self-Exponent Constraint]
For $p_i \mid n$ with exponent $a_i$, the ``self-contribution'' $2a_i - 1 \equiv 0 \pmod{3}$ requires $a_i \equiv 2 \pmod{3}$, i.e., $a_i \in \{2, 5, 8, \ldots\}$. Deviations from this are permitted only if cross-terms from $(p_j - 1)$ for other primes $p_j$ compensate.
\end{lemma}

\begin{proof}
$2a_i - 1 \equiv 0 \pmod{3} \iff 2a_i \equiv 1 \pmod{3} \iff a_i \equiv 2 \pmod{3}$, since $2^{-1} \equiv 2 \pmod{3}$.
\end{proof}

\section{Algorithm}

\begin{verbatim}
function solve(limit = 10^10):
    primes = sieve(sqrt(limit))
    result = 0
    DFS over prime factorizations of n:
        - enumerate primes in increasing order
        - for each prime p, try exponents a = 1, 2, 3, ...
        - track residues mod 3 of all prime exponents in phi(n^2)
        - prune when current_n * p^2 > limit
        - add n to result when all residues are 0 mod 3
    return result
\end{verbatim}

\section{Complexity Analysis}
\begin{itemize}
    \item \textbf{Time:} DFS with aggressive pruning; in practice $O(10^5)$ nodes. Per node, factoring $p-1$ costs $O(\sqrt{p})$.
    \item \textbf{Space:} $O(\pi(\sqrt{N}))$ for primes plus $O(k)$ stack depth ($k \le 10$).
\end{itemize}

\section{Answer}

\[
\boxed{5943040885644}
\]

\end{document}
